\chapter{Testes}\label{cap:testes}

Este capítulo registra como o comportamento do Nexo é verificado. Como nos
capítulos anteriores, o compromisso é com a honestidade: não se prescreve uma
estratégia ideal, mas se descreve o estado da prática, o que cada tipo de teste
é capaz --- e incapaz --- de detectar, e as lacunas assumidas explicitamente na
Seção~\ref{sec:test-gaps}.

O foco recai sobre a verificação automatizada. A validação com usuários reais
não foi realizada neste trabalho: o envolvimento de participantes em uma
aplicação de saúde mental exigiria aprovação prévia de comitê de ética em
pesquisa, levantamento de pessoas com transtorno de forma a aplicar com elas,
cujo todo o trâmite não coube no cronograma do projeto. Essa ausência de
avaliação externa é retomada como limitação na Seção~\ref{sec:test-gaps}.

Os dados quantitativos citados ao longo do capítulo referem-se ao estado do
repositório em meados de outubro de 2026.

\section{Vocabulário}\label{sec:test-vocab}

Quatro termos delimitam o alcance dos testes descritos adiante: unitário,
integração, \textit{job} e \textit{end-to-end} (E2E). Como os nomes e as
fronteiras entre esses níveis variam entre autores, o Quadro~\ref{tab:test-vocab}
fixa o sentido adotado neste trabalho e o relaciona aos níveis usuais da
literatura.

\begin{table}[H]
\centering
\caption{Vocabulário de testes adotado neste trabalho}
\label{tab:test-vocab}
\small
\begin{tabular}{|p{2.8cm}|p{6.6cm}|p{4.2cm}|}
\hline
\textbf{Termo} &
\textbf{O que exercita} &
\textbf{Nível usual na literatura} \\

\hline

Teste unitário &
Uma única função ou módulo, sem HTTP e sem banco de dados. &
Teste de unidade. \\

\hline

Teste de integração &
A aplicação HTTP real (rotas, validação, autenticação, controladores
e serviços) sobre um banco de dados real; apenas o que sai do processo
é simulado &
Teste de integração, em nível de API. \\

\hline

Teste de \textit{job} &
Um processador de tarefa em segundo plano, chamado diretamente com uma carga
fabricada, sobre o banco real. &
Teste de integração de componente assíncrono \\

\hline

Teste \textit{end-to-end} (E2E) &
O aplicativo empacotado, em um dispositivo, ao longo de uma jornada real de
usuário. &
Teste de sistema. \\

\hline
\end{tabular}
\source{Elaborado pelo autor (2026)}
\end{table}

\section{Estratégia de Verificação}\label{sec:test-strategy}

O código-fonte do Nexo vive em um único repositório com dois projetos
independentes: a API e o aplicativo móvel. O esforço de teste entre eles é
deliberadamente assimétrico. A API possui uma suíte automatizada e é protegida
por um portão de integração contínua; o aplicativo móvel é verificado
majoritariamente de forma manual, com uma suíte \textit{end-to-end} fina por
cima. Essa assimetria não é um descuido, e se justifica por duas prioridades.
A primeira é a capacidade de trabalho limitada: manter uma suíte automatizada
completa nas duas esferas excederia o disponível, e concentra-se o esforço onde
o retorno é maior. A segunda é que a API concentra a validação, a autenticação,
as regras de negócio e a persistência --- isto é, a garantia sobre a integridade
dos dados registrados pelo usuário.

A Tabela~\ref{tab:test-tooling} resume as ferramentas envolvidas e o ponto em que
cada uma atua.

\begin{table}[H]
\centering
\caption{Ferramentas de verificação}
\label{tab:test-tooling}
\small
\begin{tabular}{|p{4.4cm}|p{4.0cm}|p{5.0cm}|}
\hline
\textbf{Finalidade} &
\textbf{Ferramenta} &
\textbf{Onde atua} \\

\hline

Execução de testes e asserções da API &
Vitest &
API \\

\hline

Requisições HTTP nos testes &
Supertest &
Dentro dos arquivos do Vitest \\

\hline

Cobertura de código da API &
Vitest (\textit{coverage}) &
API, sem limiar mínimo \\

\hline

Banco de dados de teste &
PostgreSQL &
\textit{Container} local durante desenvolvimento; serviço efêmero no CI \\

\hline

Testes \textit{end-to-end} móveis &
Maestro &
Dispositivo Android físico, via \textit{adb} \\

\hline

\textit{Lint} do aplicativo &
ESLint (\texttt{expo lint}) &
Manual, fora do CI \\

\hline

Integração contínua &
GitHub Actions &
API \\

\hline
\end{tabular}
\source{Elaborado pelo autor (2026)}
\end{table}

\section{Suíte de Testes da API}\label{sec:test-api}

A suíte automatizada concentra-se na API e cobre três frentes: os testes
unitários, os testes de integração das rotas HTTP e os testes dos processadores
de tarefas em segundo plano. No momento da redação, o diretório
\texttt{api/tests} reunia 48 arquivos --- 17 unitários, 26 de rotas e cinco de
\textit{jobs}. A execução e as asserções ficam a cargo do Vitest; as requisições
HTTP são feitas pelo Supertest, que exercita a aplicação Express diretamente,
sem abrir uma porta de rede, porém através de sua pilha HTTP real.

\subsection{Organização e ferramentas}\label{sec:test-api-org}

A configuração do Vitest (\texttt{api/vitest.config.ts}) divide a suíte em dois
``projetos''. O projeto \texttt{unit} seleciona os arquivos com o sufixo
\texttt{.unit.test.ts} e não carrega configuração global alguma, de modo que
esses testes dependem apenas das funções que exercitam. O projeto \texttt{db}
seleciona os demais arquivos --- integração e \textit{job} ---, aplica um arquivo
de configuração global, fixa um limite de tempo por teste e faz com que esses
testes trunquem e repovoem o banco.

A cobertura de código é apenas mensurada, não imposta: o relatório
(\texttt{npm run test:coverage}) não possui limiar mínimo nem bloqueia a
integração contínua. A verificação estática de tipos (\texttt{npm run typecheck})
também roda em caráter consultivo no CI --- uma falha é registrada, mas não
interrompe o \textit{deploy}.

\subsection{Execução e banco de dados de teste}\label{sec:test-api-db}

Os testes que dependem do banco rodam contra uma instância real do PostgreSQL
--- e não contra um substituto em memória ---, em uma base dedicada de nome
sufixado com \texttt{\_test}, criada por um pequeno script e materializada com
\texttt{prisma db push}. O isolamento entre os casos é obtido por truncamento:
antes de cada teste, uma rotina executa \texttt{TRUNCATE users CASCADE},
removendo a tabela de usuários e, por cascata, todas as dependentes. Não há
\textit{rollback} transacional nem banco efêmero por execução.

O carregamento de dados é feito por funções auxiliares escritas à mão, e não por
uma biblioteca de fábricas: um módulo compartilhado
(\texttt{api/tests/test-helper.ts}) oferece as primitivas centrais --- cliente
Prisma, limpeza da base, criação de usuário e geração de token ---, e cada
arquivo define os \textit{seeders} de que necessita.

A regra para o uso de dublês na API é: simula-se apenas o que deixa o processo;
a aplicação e o banco permanecem reais. A camada de filas (BullMQ) é substituída
globalmente por um substituto que registra os \textit{jobs} enfileirados,
permitindo afirmar tanto a resposta HTTP quanto o trabalho em segundo plano
agendado. O envio de arquivos e o cliente de armazenamento são substituídos por
implementações controladas, e o disparo de e-mail e de notificações
\textit{push} é simulado pontualmente, nos arquivos que os exercitam.

Duas fronteiras permanecem deliberadamente reais: o acesso ao banco, via
Prisma, nunca é simulado; tampouco a autenticação. Os testes geram tokens JWT
válidos, assinados com o segredo de teste, e os enviam no cabeçalho de
autorização, de modo que o \textit{middleware} de autenticação real seja
executado. Essa escolha reduz o risco de um teste passar porém se comportar de
maneira distinta em produção.

\subsection{Tipos de testes}\label{sec:test-api-tipos}

A suíte da API divide-se em três grupos, retomando o vocabulário do
Quadro~\ref{tab:test-vocab}: unitários, integração e tarefas assíncronas. Cada
grupo é descrito a seguir com o que exercita, o que detecta e o que escapa ao
seu alcance.

\subsubsection{Testes unitários}

Os testes unitários exercitam uma única função ou módulo, sem HTTP e sem banco
de dados. Concentram-se na lógica pura do sistema, em grande parte extraída
deliberadamente de seus envoltórios de persistência para se tornar testável: os
agregadores dos indicadores, os construtores dos relatórios e auxiliares
compartilhados, como paginação e conversão de agendamento. Também cobrem o
gerador de código de ativação, a validade de tokens, os serializadores e as
regras de validação. Por serem rápidos e determinísticos, são o instrumento
natural para fixar casos-limite --- fronteiras de janelas de tempo, margens de
classificação e aritmética de paginação --- e servem como documentação
executável dessas regras.

A limitação desse nível é estrutural: ele só alcança a lógica isolada em uma
costura testável. Sem o banco e a camada HTTP, a maior parte do comportamento
permanece fora de alcance, e o nível pouco contribui, isoladamente, para a
confiança geral na aplicação.

\subsubsection{Testes de integração}

Os testes de integração formam a maior parte da suíte. Cada um constrói a
aplicação real, conecta um cliente Prisma real e a exercita via Supertest,
afirmando tanto a resposta HTTP quanto o estado persistido. Eles percorrem o
caminho completo de uma requisição --- roteamento, validação, autenticação,
controladores, serviços e persistência --- e, por operarem sobre um PostgreSQL
real, não sofrem o descolamento que afeta os substitutos. Afirmar a resposta e
as linhas persistidas simultaneamente verifica o efeito da operação, e não
apenas o valor de retorno.

Em contrapartida, exigem infraestrutura em execução, são mais lentos e, por
compartilharem a mesma base, rodam em série. O isolamento por truncamento é
suficiente, porém menos rigoroso do que um \textit{rollback} transacional, pois
um teste que falha no meio pode deixar resíduos.

\subsubsection{Testes de tarefas assíncronas}

Os testes de \textit{jobs} importam um processador do BullMQ e o chamam
diretamente, com uma carga fabricada, contra o banco real. Cobrem o código que
executa fora do ciclo requisição--resposta --- os \textit{workers} que geram
exportações, enviam mensagens e recalculam indicadores ---, chegando, no caso
da exportação, a descompactar o anexo produzido e inspecionar seu conteúdo.

Como o processador é invocado diretamente, a maquinaria da fila em si ---
agendamento, repetição, tentativas e ordenação --- permanece fora do alcance. A
carga é fabricada à mão e, por isso, pode divergir daquela que o enfileirador
realmente envia, sem que nada verifique uma contra a outra.

\section{Verificação do Aplicativo Móvel}\label{sec:test-mobile}

Em linha com a assimetria descrita na Seção~\ref{sec:test-strategy}, o
aplicativo móvel passou a contar com testes unitários em estágio inicial, ainda
em andamento e sem a cobertura da suíte da API. A verificação cotidiana
permanece predominantemente manual, em um dispositivo físico, contra a pilha
local em execução --- o aplicativo roda como \textit{build} de desenvolvimento
do Expo, com dependências nativas do Android. Essa verificação captura o que os
testes automatizados raramente alcançam --- problemas visuais, de leiaute e de
desempenho no dispositivo ---, mas não é repetível, não protege contra
regressões e permanece invisível à integração contínua.

\subsection{Suíte \textit{end-to-end} (Maestro)}\label{sec:test-mobile-e2e}

Sobre a verificação manual há uma suíte \textit{end-to-end} opcional, baseada
no Maestro, ferramenta declarativa de automação de interface para dispositivos
móveis. A suíte cobre uma única jornada central --- login, registro de humor e
saída --- por três fluxos em YAML, marcados como \texttt{smoke} e executados
por um pequeno \textit{harness} (\texttt{frontend/scripts/e2e.sh}), exposto
como \texttt{npm run test:e2e}. Os fluxos são dirigidos por identificadores de
teste (\texttt{testID}) inseridos no código, e não pelo texto visível, o que os
torna mais robustos a mudanças de interface.

A execução depende de um dispositivo Android físico conectado por \texttt{adb},
de um \textit{build} de desenvolvimento, da API em funcionamento com um usuário
de teste semeado e do Metro servindo o aplicativo; não há caminho sem cabeça
nem integração ao CI, de modo que a suíte é executada manualmente e com pouca
frequência. Cada fluxo limpa o estado do aplicativo e recarrega o
\textit{bundle}, tornando a execução lenta, e a cobertura permanece fina: um
único caminho feliz, sem casos de erro. A suíte não está integralmente verde no
momento da redação: os fluxos de login e de saída passam, mas o de registro de
humor falha ao adicionar um componente de humor, porque a tela correspondente é
abandonada pelo aplicativo pouco depois de ser aberta --- causa ainda não
identificada.

\section{Integração Contínua}\label{sec:test-ci}

A integração contínua do Nexo \cite{Humble_Farley_2011} é executada no GitHub
Actions e testa somente a API, coerente com a estratégia assimétrica. O fluxo
de implantação (\texttt{.github/workflows/deploy.yml}) dispara apenas quando há
mudanças em \texttt{api}: sobe um PostgreSQL efêmero, instala as dependências,
prepara a base de teste, executa a verificação de tipos em caráter consultivo e
roda a suíte com o Vitest. A construção da imagem e a implantação em produção
são condicionadas ao sucesso dessa etapa, de modo que uma suíte reprovada
impede a publicação, conforme o Capítulo~\ref{sec:deploy-and-ci}. O aplicativo
móvel permanece fora da integração contínua: \textit{lint} (\texttt{expo lint})
e verificação de tipos (\texttt{tsc --noEmit}) são executados apenas
manualmente, e a suíte \textit{end-to-end} não é disparada por fluxo algum.

\section{Cobertura e Lacunas}\label{sec:test-gaps}

Em síntese, a verificação automatizada cobre a superfície HTTP e a persistência
da API, a sua lógica pura e os seus trabalhadores em segundo plano; a
verificação do aplicativo móvel restringe-se, por ora, a uma jornada manual de
fumaça e a testes unitários ainda em consolidação. Permanecem pouco cobertas a
lógica do \textit{frontend} entre o nível de unidade e o \textit{end-to-end}, a
maquinaria da fila em si (agendamento, repetições e ordenação) e os caminhos de
erro do aplicativo.

Persistem, ainda, três limitações que merecem registro. A primeira é a ausência
de validação com usuários reais, motivada pelo trâmite de comitê de ética em
pesquisa descrito na abertura deste capítulo: as especificações automatizadas
atestam o comportamento programado, mas não a adequação do sistema às
expectativas do público-alvo. A segunda é a cobertura de código medida, mas não
imposta, o que orienta o trabalho sem jamais reprovar uma \textit{build} ou
\textit{deploy}. A terceira são os riscos de descolamento já mencionados: os
dados de teste escritos à mão em cada arquivo e a cópia espelhada dos tipos da
API mantida no \textit{frontend} podem divergir da fonte de verdade sem que
nenhum teste acuse. Essas lacunas são assumidas de forma explícita e refletem
uma escolha de priorização, não uma garantia de ausência de defeitos, e que
podem ser corrigidas em incrementos no futuro.
